\section{Mẹo hay}

\begin{itemize}
  \item \textbf{Chèn nhanh đoạn mẫu}: Gõ \texttt{\textbackslash fig},
    \texttt{\textbackslash table}, \texttt{\textbackslash eq},
    \texttt{\textbackslash ilist} (danh sách gạch đầu dòng) hoặc
    \texttt{\textbackslash elist} (danh sách đánh số) rồi nhấn \phim{Enter}.
    Sau đó, dùng phím \phim{Tab} để nhảy nhanh qua các vị trí cần điền; mọi gói
    thư viện cần thiết đều sẽ được tự động thêm vào.
  \item \textbf{Tìm kiếm trên toàn dự án}: Nhấn \phim{Ctrl}+\phim{Shift}+\phim{F}.
  \item \textbf{Tìm hình ở tệp khác}: Nhấn \phim{Ctrl}+\phim{P}, sau đó gõ
    \texttt{!fig} kèm vài từ khóa có trong chú thích của hình.
  \item \textbf{Đồng bộ giữa mã và PDF}: Nhấn \phim{Ctrl}+\phim{Alt}+\phim{J}
    để cuộn bản PDF tới đúng vị trí con trỏ trong mã; ngược lại, nhấp đúp vào
    một chữ trên PDF để nhảy về dòng mã tương ứng.
  \item \textbf{Xem lịch sử phiên bản}: Nhấn \phim{Ctrl}+\phim{Shift}+\phim{H}
    để xem lại các phiên bản đã lưu và khôi phục khi cần thiết.
  \item \textbf{Cộng tác trực tuyến}: Bấm nút \nut{Share} để mời người khác
    cùng viết. Bạn có thể thấy vị trí con trỏ của nhau và trao đổi trực tiếp ở
    thanh bên phải.
  \item \textbf{Đếm số từ}: Tổng số từ của tài liệu luôn được cập nhật ở góc
    dưới bên phải màn hình.
\end{itemize}

\begin{chuy}
SumiTeX tự động lưu mọi nội dung bạn gõ (trạng thái \nut{Saved} hiển thị ở góc
dưới bên phải). Do đó, thao tác biên dịch chỉ dùng để cập nhật bản PDF mới chứ
không phải để lưu bài.
\end{chuy}